% Folder 136, batch 17 — archivists' pages 321–340.
% Pass: Opus 5.5 (claude-opus-5-5), 2026-10-03 — first pass, unchecked against the pages by a human.
%
% Page 335 carries only the word « Gribouillis » in his hand: treated as a
% cover sheet (no \page), the word becomes the section heading before p. 336.
% Page 340 was scanned upside down and is read turned.
\documentclass[11pt,a4paper]{article}
\input{../preamble/grothendieck}

\folder{136}
\batch{17}
\pages{321}{340}
\foldertitle{Complexe de De Rham à puissance divisée [conférence de 1976 à l’IHÉS] : notes manuscrites (s.d.).}
\dating{[à partir de 1975-1976]}
\watermark{Édition de démonstration}

\begin{document}

\page{321}
\note{la page continue un calcul commencé avant le lot : la notation $F_A$, $u_p$, $\Gamma^p$ est celle des pages précédentes.}

\struck{$\Gamma^p(f)\circ u_p$} \quad \struck{$\Gamma^p(g)$}

\[
\Gamma^p(F_A(X))\otimes\Gamma^q(F_A(X)) \longrightarrow F_A\Gamma^p(X)\otimes F_A\Gamma^q(X)
\]
\[
\Gamma^{p+q}F_A(X) \longrightarrow F_A\Gamma^{p+q}F(X)
\]
\note{les deux lignes sont reliées par des flèches verticales ; le second membre de la seconde est surchargé et peu sûr.}

$f, g : A \rightrightarrows X$,
\[
(\Gamma^p f\circ u_p)\otimes(\Gamma^q g\circ u_q), \qquad
\bigl[((\Gamma^p f)\circ u_p)\otimes(\Gamma^q g\circ u_q)\bigr]\circ\Delta
\]
$A\to\Gamma^pX$, \quad $A\to\Gamma^qX$, \quad $f^{(p)}\otimes g^{(p)}$\note{sic, $g^{(p)}$ pour $g^{(q)}$.}

\begin{tikzcd}[column sep=small, row sep=small, nodes={font=\scriptsize}]
A \arrow[r, "\Delta"] \arrow[d] & A\otimes A \arrow[r, "{(\Gamma^p f\circ u_p)\otimes(\Gamma^q g\circ u_q)}"] & \Gamma^p(X)\otimes\Gamma^q(X) \arrow[d] \\
{} & {} & \Gamma^{p+q}(X)
\end{tikzcd}

Une flèche courbe de $A$ vers $\Gamma^{p+q}(X)$ porte : « définit $f^{(p)}g^{(q)} \in \Gamma^{p+q}F_A(X)$ ».

\begin{tikzcd}[column sep=small, row sep=small, nodes={font=\scriptsize}]
\Gamma^p(F_A(X))\otimes\Gamma^q(F_A(X)) \arrow[r] \arrow[d] & F_A\Gamma^p(X)\otimes F_A\Gamma^q(X) \arrow[d] \\
\Gamma^{p+q}F_A(X) \arrow[dr] & F_A(\Gamma^p(X)\otimes\Gamma^q(X)) \arrow[d] \\
{} & F_A(\Gamma^{p+q}(X))
\end{tikzcd}

avec $f^{(p)}\otimes g^{(q)} \mapsto \Gamma^p(f)\circ u_p\otimes\Gamma^q(g)\circ u_q$, et à gauche $f^{(p)}g^{(q)}$.

Il suffit de \uncertain{le voir} \ill{} $X = Y = A$, $f = g = \mathrm{id}$ :

\begin{tikzcd}[column sep=small, row sep=small, nodes={font=\scriptsize}]
A \arrow[r, "\Delta"] & A\otimes A \arrow[r, "u_p\otimes u_q"] & \Gamma^p(A)\otimes\Gamma^q(A) \arrow[d] \\
{} & {} & \Gamma^{p+q}(A)
\end{tikzcd}

la flèche courbe de $A$ vers $\Gamma^{p+q}(A)$ « définit $(\mathrm{id}_A)^{(p)}(\mathrm{id}_A)^{(q)} \in \Gamma^{p+q}F_A(A)$ »
\struck{i.e. $\beta_{pq}(\mathrm{id}_A^{(p+q)})$, donc $c_{p,q}\,u_{p+q}$}.
\[
1^{(p)}1^{(q)} = c_{pq}\,1^{(p+q)}
\]
\marginal{double trait vertical dans la marge gauche, en regard de la ligne suivante}
\[
\underbrace{\Gamma^*F(A)}_{\otimes\text{-anneau}} \longrightarrow \underbrace{F(\Gamma^*(X))}_{\otimes\text{-anneau}}
\]
car $F$ commute à $\otimes$.

\page{322}
(1) $\Longleftrightarrow$ (3) \uncertain{déjà vu}.

(2) $\Longleftrightarrow$ (3). La donnée d'\uncertain{homomorphismes} fonctoriels
\[
F_A \xrightarrow{u_p} G = F_A\circ\Gamma^p
\]
équivaut à la donnée d'un élément
\[
u_p \in G(A) = \mathrm{Hom}(A,\Gamma^p(A)).
\]
Je \uncertain{laisse} \ill{} \uncertain{la compatibilité} \ill{} de (2). Les conditions $f^{(0)} = 1$, $f^{(1)} = f$ correspondent à $u_0 = 1$, $u_1 = \mathrm{id}$. La relation $(f+g)^{(p)} = \sum f^{(i)}g^{(j)}$ : il suffit de voir le cas universel
\[
M = A\times A, \quad f = \mathrm{pr}_1, \quad g = \mathrm{pr}_2,
\]
$f+g = \Delta : A \to A\times A = M$,
\[
(f+g)^{(p)} = \Gamma^p(\Delta)\circ u_p \;:\; \Gamma^p(A) \longrightarrow \Gamma^p(M)
\]
\note{des surcharges barrées dans les trois dernières formules ne sont pas lisibles.}

\begin{tikzcd}[column sep=small, row sep=small, nodes={font=\scriptsize}]
A \arrow[r, "u_p"] \arrow[d, "\Delta_A^{(p)}"'] & \Gamma^p(A) \arrow[r, "\Gamma^p(\Delta)"] & \Gamma^p(A\times A) \arrow[d, "\simeq"] \\
A\times\cdots\times A \arrow[rr] & {} & \sum_{i+j=p}\Gamma^iA\otimes\Gamma^jA
\end{tikzcd}

\note{le produit $A\times\cdots\times A$ porte « $p$ » sous une accolade ; une flèche oblique entre parenthèses est marquée $\Delta_{\Gamma^*A}$ ; sous la ligne du bas : « $\mathrm{Gr}\; A \to \Gamma^i(A)\otimes\Gamma^j(A)$ ».}

On trouve
\[
\Gamma^p(\Delta)(u_p) = \sum_{i+j=p} u_i\otimes u_j .
\]

\page{323}
Deux questions de p.d. (\uncertain{anneaux commutatifs}).

(1) Soient $A$ une cogèbre \uncertain{associative} \uncertain{commutative}, $(B,J)$ une algèbre à idéal à p.d. (p.ex. $(\Gamma^*M,\Gamma^+M)$). On sait que $\mathrm{Hom}_k(A,M)$ a une structure d'anneau --- \uncertain{même} de $k$-algèbre, et $\mathrm{Hom}_k(A,J)$ \uncertain{en est} un idéal. (Si $A$ est libre de type fini, on trouve l'anneau $\check A\otimes B$ avec l'idéal $\check A\otimes J$, qui a des p.d.) Cet idéal a-t-il des p.d.~?

(2) Prenons $B = \Gamma^*A$, $J = \Gamma^+A$, et soit $u_1 \in \mathrm{Hom}_k(A,B)$ l'\uncertain{application identique} $A\to A$. Considérons les $u_n = u_1^{(n)}$. Satisfont-ils
\[
\gamma^A_{pq}\circ u_{pq} = \Gamma^p(u_q)\circ u_p \;?
\]
Sont-ils (pour \uncertain{$n$, $p\ge0$}) les seuls \uncertain{éléments} satisfaisant cette relation, plus $u_pu_q = c_{pq}u_{p+q}$, et pour lesquels $u_1$ est comme dessus~?

(3) Soit $A$ une \uncertain{algèbre} \uncertain{augmentée} (\uncertain{ass., comm.}), $B$ une cogèbre à coproduits \uncertain{divisés} pour la \uncertain{coloi} de $B'\to B$ (\uncertain{noyau} de \uncertain{coaugm.}), p.ex. $B = \mathrm{Sym}^*M$, \uncertain{prenons} $\mathrm{Sym}^*(0) = k$ … Donc sur $\mathrm{Hom}_k(B,A)$ il y a une structure \struck{de p.d. relativement} d'algèbre, à p.d. sur
\[
\mathrm{Ker}\,\bigl(\mathrm{Hom}_k(B,A)\to\mathrm{Hom}_k(B',A)\bigr).
\]

\page{324}
3) Donnée d'une application (pour tt $p\ge0$)
\[
\mathrm{Hom}_k(A,M) \xrightarrow{u_p} \mathrm{Hom}_k(A,\Gamma^pM)
\]
(\add{non} linéaire) fonctorielle en $M$, notée \uncertain{aussi} \struck{faisons} \ill{}
\[
f \longmapsto f^{(p)}
\]
et satisfaisant les relations $f^{(0)} = 1$, $f^{(1)} = f$, et
\[
\begin{cases}
(f+g)^{(p)} = \sum_{i+j=p} f^{(i)}g^{(j)} \\
(\lambda f)^{(p)} = \lambda^p f^{(p)} \quad (\text{automatique~!})
\end{cases}
\]
et …

4) Donnée d'une loi de p.d. sur
\[
\mathrm{Hom}_k(A,\Gamma^+M) \subset \mathrm{Hom}_k(A,M)
\]
$\forall M$, fonctorielle en $M$, et telle que …

5) Loi à p.d. sur \struck{$(A\otimes\Gamma^+A)$} \add{$\mathrm{Hom}(A,\Gamma^+A)$} $\subset \mathrm{Hom}(A,\Gamma^*A)$ (\uncertain{ann. commutatif}).

6) Donnée, pour \uncertain{toute} $k$-algèbre $B$ (avec idéal $J$ à p.d.), d'une loi de p.d. sur
\[
\mathrm{Hom}_k(A,J) \subset \mathrm{Hom}_k(A,B),
\]
fonctorielle en $(B,J)$ ……
\marginal{double trait vertical dans la marge droite, en bas de page}

\page{325}
\note{feuillet à l'encre bleue, écrit en travers, fait de calculs épars ; on transcrit les formules dans l'ordre approximatif de la page. Le sujet change : groupes $Sl(2)$, classes d'Euler, connexions sur une courbe.}

\struck{$GL(2,\mathbb{R}) \to GP(\ldots)$}
\[
\begin{array}{ccccccccc}
1 & \to & \mu_2 & \to & Sl(2,\mathbb{C}) & \to & GP(1,\mathbb{C}) & \to & 1 \\
  &     & \uparrow &  & \uparrow &  & \uparrow & & \\
1 & \to & \mu_2 & \to & Sl(2,\mathbb{R}) & \to & GP(1,\mathbb{R}) & \to & 1
\end{array}
\]
$T_{\mathbb{C}}$, $c_1$ ; \quad $E_1$, \; $T^0_{\mathbb{R}} \simeq S^1 = SO(2,\mathbb{R})$, \; $c_1$ ; \quad $\tfrac12 E_1 = E'_1$, \quad $\tfrac12 c_2 = c'_2$.
\[
\begin{pmatrix} a & b \\ -b & a \end{pmatrix}, \quad a^2+b^2 = 1, \qquad
(a+ib)(a-ib) = 1, \quad \lambda\mu = 1,
\]
\[
a = \frac{\lambda+\mu}{2} = \frac12(\lambda+\lambda^{-1}), \qquad
b = \frac{\lambda-\mu}{2i} = \frac1{2i}(\lambda-\lambda^{-1}).
\]
\[
\Omega\otimes\Omega, \qquad \mathbb{P}(\Omega+\Omega^{-1}), \qquad
\Omega^{-1}\;\;\Omega^{-2} \,/\, \Omega,\;\mathcal{O}, \qquad
\mathbb{P}^3, \qquad H^1(X,\Omega) \simeq \mathbb{C}
\]
\[
\cdots \to \Omega^{\otimes n} \to P^n_{X/k} \to P^{n-1}_{X/k} \to 0
\]
$Gl \to GP$, \quad $P \to P'$, \quad $\check E$, $\mathcal{O}_X$ ; \quad (n) épi $E \to \mathcal{O}_X$ ;
\[
E \longrightarrow E\otimes\Omega^1_{X/S}, \qquad E \to \mathcal{O}_X, \qquad
\mathrm{Sym}^n\Omega^1 \to P^n_{X/k} \xrightarrow{D} L,
\]
\[
a_nf^{(n)} + \cdots + a_1 f^{(n-1)} \ldots, \qquad Df + \alpha f = 0, \qquad \mathcal{O}_X + P^{n+}_{X/k}
\]
\ill{}
\note{plusieurs indices et la ligne $a_nf^{(n)}+\cdots$ sont incertains ; le haut de la page est barré de deux traits obliques.}

\page{326}
$X$ courbe alg.\ /$\mathbb{C}$ \uncertain{projective lisse connexe} de genre $g\ge2$. Il y a sur $X$ un \struck{fibré principal} \add{torseur} canonique $P$ de groupe $GP(1)_{\mathbb{C}} = G'_{\mathbb{C}}$ à connexion intégrable (\uncertain{équivalente} à structure de torseur), obtenu ainsi : $\forall x\in X$, $\tilde X(x)$ \uncertain{rev.} universel de $X$ en $x$, $\tilde X(x) \Longleftrightarrow$ \struck{\ill{}} \add{demi-plan de} Poincaré $D$ ; \uncertain{l'ensemble} des \uncertain{isos} $\mathrm{Iso}(D_0,(\tilde X(x),\tilde x))$ est un torseur à droite sous $\mathrm{Aut}(D_0) \simeq$ \struck{$GP(1)_{\mathbb{R}}$} $G'_{\mathbb{R}}$, d'où par \uncertain{ext.} un torseur à droite sous $G'$, \add{soit $P_x$} \struck{\ill{}}. On voit que les $P_x$ ($x\in X(\mathbb{C})$) sont les fibres d'un fibré principal \struck{\ill{}} \add{\uncertain{torsoriel}} à structure plate (celui associé : $\pi_1 \to G'_{\mathbb{R}} \subset G'_{\mathbb{C}}$).

\drawing{un disque hachuré marqué d'un point $\tilde x$, dans un carré, avec une boucle qui en sort}

N.B. $\forall x\in X$, $\exists$ une et une seule \ill{} \uncertain{application analytique} de $\tilde X(x)$ dans $T_x(X)$ dont a) l'image soit un disque, b) l'application tangente en $\tilde x$ soit l'identité.

\[
\pi_1(X,x) \subset G' \subset G'_{\mathbb{C}}, \qquad
\begin{pmatrix} a & b \\ c & d \end{pmatrix}.z = , \qquad
\begin{pmatrix} a & b \\ -b & a \end{pmatrix}
\]
\struck{$\deg$} $T(x,q)$,
\[
P(x,\rho q) = P(x,q)\wedge^{i_x(\rho)} G'_{\mathbb{C}}, \qquad \rho \in S^1_{\mathbb{C}} \subset G'_{\mathbb{C}},
\]
\[
\begin{array}{ccc}
S_{\mathbb{R}} & \subset & G'_{\mathbb{R}} \\
\cap & & \cap \\
S_{\mathbb{C}} & \subset & G'_{\mathbb{C}}
\end{array}
\qquad Q_{\mathbb{R}},\; P_{\mathbb{R}},\; Q_{\mathbb{C}},\; P_{\mathbb{C}}, \qquad
\mathbb{C}^* \subset Gl(2,\mathbb{R}) \to Gl(2,\mathbb{C})
\]
\note{la dernière inclusion est mal lue ; le texte porte aussi $P(x,q) \simeq P(x,x')$ et une flèche $G'_{\mathbb{C}} \to G'_{\mathbb{C}}$ marquée $i(\ldots)$.}

$E \to E\otimes\Omega$ \;(deux fois).

a) fibré inversible : $\Omega^{-1}$ \qquad $G_m \to GP_{\mathbb{C}}$

b) fibré $\mathbb{P}(\Omega^{-1}+\Omega)$ avec connexion intégrable.

\[
\begin{pmatrix} \Omega & 0 \\ \Omega^3 & \Omega \end{pmatrix}
\]
Encadré : Trouver une connexion intégrable (\uncertain{alg.}) canonique sur
\[
\mathbb{P}(\Omega^{-1}+\Omega) = \overline{V(\Omega^2)}
\]
\[
\Omega^{-1}\oplus\Omega \longrightarrow (\Omega^{-1}\oplus\Omega)\otimes\Omega = \Omega^0_X + \Omega^2
\]
$(D,\omega) =$ \ill{} $(\varpi,\Delta)$

\page{327}
\[
\begin{array}{ccc}
\tilde G_d & \to & \tilde G \\
\vdots & & \vdots \\
\mathfrak{z}_d & \to & \mathfrak{z} \quad \mathbb{Z}/2\mathbb{Z} \\
G_d & \to & G = Sl(2,\mathbb{R}) \\
\downarrow & & \downarrow \\
G'_d & \to & G' = GP(1,\mathbb{R}) \\
\downarrow & & \downarrow \\
1 & & 1
\end{array}
\qquad
\begin{array}{c}
\mathbb{Z} \\ \downarrow \\ \mathfrak{z} = \mathbb{Z}/2\mathbb{Z}
\end{array}
\]
\note{lecture incertaine de $G' = GP(1,\mathbb{R})$. $\mathfrak{z}$ rend la lettre bouclée en forme de 3 par laquelle il note le noyau $\mathbb{Z}/2\mathbb{Z}$.}
\[
H^4(B_G,\mathbb{Z}) \simeq \mathbb{Z} \ni W, \qquad
\uparrow \qquad \uparrow 2, \qquad
H^2(B_{G'},\mathbb{Z}) \simeq \mathbb{Z} \ni \tfrac12 W
\]
\struck{$H^2(B_G,\mathbb{T})$} \quad $H^4(B_G,\mathfrak{z})$, \; $\mathbb{Z}/2.2 = 0$
\[
H^2(B_{G'},\mathfrak{z}) \simeq \mathfrak{z} \ni \xi
\]
\[
\begin{array}{ccc}
H^2(B_{G'},\mathfrak{z}) & & \\
\downarrow & & \\
H^2(B_{G'_d},\mathfrak{z}) & \xleftarrow{\;\sim\;} & H^2(B_{G'_d},\mathfrak{z}_d) \\
\downarrow & & \\
H^2(\pi,\mathfrak{z}) & & H^2(X,\mathfrak{z})
\end{array}
\]
\struck{$\tfrac12 W = 2g-2$}

\[
\begin{array}{ccccc}
\mathbb{Z} \simeq H^2(X,\mathbb{Z}) & \leftarrow & H^2(B_{G'},\mathbb{Z}) & \ni & W' \\
\downarrow & & \downarrow & & \\
\mathbb{Z}/2\mathbb{Z} \quad H^2(X,\mathfrak{z}) & \leftarrow & H^2(B_{G'},\mathfrak{z}) \simeq \mathbb{Z}/2\mathbb{Z} & \ni & y \\
\updownarrow 2 & & \downarrow & & \\
H^2(\pi,\mathfrak{z}) & \leftarrow & H^2(B_{G'_d},\mathfrak{z}) & &
\end{array}
\]
\note{le carré du haut porte « comm. » ; la flèche $\updownarrow 2$ est telle qu'écrite.}

À l'encre bleue :
\[
1 \to \mathfrak{z} \to G \to G' \to 1, \qquad H^2(B_{G'},\mathfrak{z})
\]
\[
X \to B_{G'_d} \to B_{G'}, \qquad H^2(B_{G'}) \ni E' \simeq H^2(B_{T'})
\]
\[
X \to B_\pi \to B_{G'} \leftarrow B_{T'}, \qquad
\pi \hookrightarrow GP = G' \to S^1 = T'
\]
\[
\uparrow \qquad \uparrow \qquad Sl = G \to S^1 = T, \qquad E' \mapsto 2g-2 \sim g-1
\]
\[
\mathbb{Z} \simeq H^2(X,\mathbb{Z}) \leftarrow H^2(\pi,\mathbb{Z}) \leftarrow H^2(B_{G'},\mathbb{Z}) \ni E'
\]
\[
\begin{array}{ccccccc}
X & \to & B_\pi & \to & B_{G'} & \leftarrow & B_{T'} \\
 & & & & \uparrow & & \uparrow \\
 & & & & B_G & \leftarrow & B_T
\end{array}
\qquad
\begin{array}{c}
H^2(B_{T'},\mathbb{Z}) \simeq \mathbb{Z} \ni E' = \tfrac12 E \\
\downarrow \\
H^2(B_T,\mathbb{Z}) \simeq \mathbb{Z} \ni E
\end{array}
\]
\note{en marge à gauche : $2g-2 \sim g-1$.}

\page{328}
\note{retour au sujet des pp.~321--324, à l'encre noire.}

\begin{tikzcd}[column sep=small, row sep=small, nodes={font=\scriptsize}]
A \arrow[r, "u_{pq}"] \arrow[d, "u_p"'] & \Gamma^{pq}(A) \arrow[d, "\gamma^A_{pq}"] \\
\Gamma^p(A) \arrow[r, "\Gamma^p(u_q)"'] & \Gamma^p(\Gamma^q(A))
\end{tikzcd}

À droite : $\Gamma^p\mathrm{Hom}(A,M) \to \mathrm{Hom}(A,\Gamma^pM)$.

\begin{tikzcd}[column sep=small, row sep=small, nodes={font=\scriptsize}]
A \arrow[r, "f"] \arrow[d, "u_p"'] & \Gamma^*M & {} \\
\Gamma^pA \arrow[r, "\Gamma^pf"'] & \Gamma^p(\Gamma^*M) \arrow[r] & \Gamma^*M
\end{tikzcd}

\begin{tikzcd}[column sep=small, row sep=small, nodes={font=\scriptsize}]
A \arrow[r, "f"] \arrow[d, "u_p"'] & J & {} \\
\Gamma^p(A) \arrow[r, "\Gamma^p(f)"'] & \Gamma^p(J) \arrow[r] & J
\end{tikzcd}

Les $A$ \uncertain{munis de} \ill{} \uncertain{cocommutative}, \ill{} par le foncteur $F_A(M) = \mathrm{Hom}_k(A,M)$,
\[
\mathrm{Ab}_k \xrightarrow{\;F\;} \mathrm{Ab}_k
\]
\ill{} d'une structure de \uncertain{quasi-commutation} \ill{} compatible avec \ill{}, \uncertain{comm.}).

\struck{Conditions} \add{Données} \uncertain{équivalentes}.

(1) Donnée d'une \uncertain{quasi-commutation} de $F$ aux foncteurs $\Gamma^p$, avec \uncertain{compatibilité} aux $\beta_{pq}$ (structure multiplicative de $\Gamma$) et aux $\gamma_{pq}$ (structure \ill{}), enfin avec condition de \uncertain{comp.} en degré 0, 1.

(2) Donnée d'une famille d'éléments
\[
u_p \in F_A(\Gamma^p(A)) \;:\; A \to \Gamma^p(A)
\]
satisfaisant
\[
\begin{cases}
u_pu_q = c_{pq}\,u_{p+q} \\
\Gamma^p(u_q)\circ u_p = \gamma^A_{pq}\circ u_{pq}
\end{cases}
\]
\ill{} d'un homomorphisme \struck{\ill{}} de \uncertain{$k$-alg.\ gr.\ divisées}
\[
k\{T\} \xrightarrow{\;u\;} \mathrm{Hom}(A,\Gamma^*(A))
\]
\note{sous $\mathrm{Hom}(A,\Gamma^*(A))$, une accolade : « \uncertain{algèbre} \uncertain{quasi-commutative} \ill{} de la \uncertain{cogèbre} $A$ et \ill{} \uncertain{gradué} de $\Gamma^*(A)$ ».}
satisfaisant
\[
\Gamma^p\bigl(u(T^{(q)})\bigr)\circ u(T^{(p)}) = \gamma^A_{pq}\,u(T^{(pq)}).
\]

\page{329}
\[
\begin{array}{c}
\Gamma^p(\Gamma^q(X)) \\
\uparrow{\scriptstyle\gamma^X_{p,q}} \\
\Gamma^{pq}(X)
\end{array}
\qquad
\begin{aligned}
\Delta &: A \to A\otimes A \\
u_p &: A \to \Gamma^pA
\end{aligned}
\]
\[
\beta^X_{p,q} : \Gamma^p(X)\otimes\Gamma^q(X) \longrightarrow \Gamma^{p+q}(X)
\]
\[
\alpha^{X,Y}_{pq} : \Gamma^p(X)\otimes\Gamma^q(Y) \xrightarrow{\;\sim\;} \Gamma^{2}(X\times Y)
\]
\note{l'exposant du but est surchargé ; on lit, sous réserve, un $2$.}
\[
F(X\otimes Y) \xleftarrow{\;\tau_{X,Y}\;} F(X)\otimes F(Y), \qquad
F(\Gamma^p(X)) \xleftarrow{\;\sigma^X_p\;} \Gamma^pF(X)
\]
« ok ok »

\note{la moitié inférieure de la page est traversée de quatre longs traits obliques et d'un trait presque vertical, qui semblent l'annuler ; elle est transcrite telle quelle.}

Soit $A$ \uncertain{algèbre} sur un anneau $k$ \struck{\ill{}}, on définit une \uncertain{flèche can.} \struck{\ill{}}
\[
\mathrm{Sym}^*A \longrightarrow A
\]
\ill{} \uncertain{restriction}
\[
\mathrm{Sym}^n(A) \xrightarrow{\;u_n\;} A \qquad (n\ge0)
\]
et on prendra
\[
u_n(x_1\cdot x_2\cdots x_n) = \frac{1}{n!}\,x_1x_2\cdots x_n
\]
\note{sous $x_1\cdot x_2\cdots x_n$ : « produit dans $\mathrm{Sym}^n(A)$ ».}
\uncertain{je dis que} \ill{} \uncertain{le diagramme}
\[
\begin{array}{ccc}
\mathrm{Sym}^{p+q}A & \xrightarrow{\;c_{pq}u_{p+q}\;} & A \\
\downarrow{\scriptstyle\Delta} & & \uparrow{\scriptstyle\text{mult.\ de }A} \\
\mathrm{Sym}^p A\otimes\mathrm{Sym}^q A & \xrightarrow{\;u_p\otimes u_q\;} & A\otimes A
\end{array}
\]
\ill{} \uncertain{une} \ill{} \uncertain{can.} de \struck{$(\Gamma^\alpha(A)\otimes A)^\wedge$}
\struck{$\Gamma^\alpha(V(A))$} $\mathrm{Hom}(\mathrm{Sym}^*A, B)$ \ill{}
\[
= \Gamma^*(\check A)\otimes_k B
\]
\ill{} $\check A\otimes_k B$ ;
$\Gamma^\alpha(\check A)\otimes B$, \uncertain{p.\,ex.} de $A \xrightarrow{h} B$ \uncertain{$k$-lin.}
$\mathrm{Sym}^*A$, \quad $\mathrm{Sym}^*($ \ldots
\note{la page s'arrête sur ce début de formule.}

\page{330}
\struck{Il suffit de \ill{}, \ill{} \ill{} \ill{} fonctorielle en $M$ : $\mathrm{Hom}(A,M) \xrightarrow{u}$}
\note{seule ligne de la page, barrée par un trait ondulé ; le reste du feuillet est blanc.}

\page{331}
\[
\Gamma^\alpha F(X) \xrightarrow{\;u_\alpha\;} F\Gamma^\alpha X, \qquad
\Gamma^pF(X) \xrightarrow{\;u_p\;} F\Gamma^p(X)
\]
\struck{$x\in$} $F(X) \to F\Gamma^p(X)$ (flèche marquée \ill{}), \quad $x\in F(X)$, \uncertain{degré} $p$
\[
u_p\bigl(\gamma^p_{F(X)}(x)\bigr) = F(\gamma^p_X)(x), \qquad
X \longrightarrow \Gamma^p(X) \;\; \text{degré } p
\]
\[
u_\alpha\bigl(\gamma^p_{F(X)}(x)\bigr) = F(\gamma^p_X)
\]

\begin{tikzcd}[column sep=small, row sep=small, nodes={font=\scriptsize}]
\Gamma^pF(X) \arrow[r, "u_p"] \arrow[d, "\xi\mapsto\xi^{(q)}"'] & F(\Gamma^p(X)) \arrow[d, "F(\eta\mapsto\eta^{(q)})"] \\
\Gamma^{pq}(F(X)) \arrow[r, "u_{pq}"'] & F(\Gamma^{pq}(X))
\end{tikzcd}

\note{à gauche, une accolade relie $\Gamma^pF(X)$ à $\Gamma^q(\Gamma^pF(X))$, avec $\xi\mapsto\gamma^q(\xi)$ ; un « \& » suit le $u_p$ en tête de ligne.}

\[
u_p(\gamma^p(x)) = F(\gamma^p_x)(x)
\]
\[
F(\gamma^{q,p}_X)\bigl(F(\gamma^p_X)(x)\bigr)
\]
\[
F(\gamma^{q,p}_X\gamma^p_X)(x)
\]
\[
\begin{aligned}
F(\alpha_{pq}\gamma^{pq}_X)(x) &= \alpha_{pq}F(\gamma^{pq}_X)(x) \\
&= \alpha_{pq}\,\ldots \\
&= u_{pq}(\alpha_{pq}\gamma^{pq}(x)) \\
&= u_{pq}\bigl(\gamma^{q,p}_{F(X)}(\gamma^p_{F(X)}(x))\bigr)
\end{aligned}
\]
avec
\[
\alpha_{pq} = \frac{(pq)!}{(p!)^q\,q!}
\]
\ill{} $=$ \ill{}
\[
u_\alpha\bigl(\gamma^{\alpha_1}(x_1)\cdots\gamma^{\alpha_n}(x_n)\bigr)
\]
\[
u(\lambda)^{pq}\,u\bigl(\alpha_{pq}x^{(pq)}\bigr), \qquad
\lambda^{pq}\alpha_{pq}\,u(\ldots)^{(pq)}, \qquad
u(x)^{\ldots}\,\alpha_{pq}\,u(x)^{[pq]}
\]
\note{dans la deuxième expression, un argument barré illisible ; l'exposant de $u(x)$ dans la troisième est illisible. Ces trois dernières expressions, au bas de la page, sont reliées par des signes $=$ verticaux ; leur lecture est très incertaine.}

\page{332}
\[
u_\alpha\Bigl(\bigl(\textstyle\sum\xi_i\bigr)^{(q)}\Bigr) \overset{?}{=} \Bigl(\textstyle\sum u_\alpha(\xi_i)\Bigr)^{(q)}
\]
\[
u\Bigl(\textstyle\sum \xi_1^{(q_1)}\cdots\xi_r^{(q_r)}\Bigr) = \textstyle\sum u(\xi_1)^{(q_1)}\cdots u(\xi_r)^{(q_r)}
\]
On se ramène à prouver que
\[
u_\alpha\bigl(\xi_i^{(q_i)}\bigr) = u_\alpha(\xi_i)^{(q_i)}
\]
pour les \uncertain{monômes}, \ill{} $q$ ; \uncertain{tels} \uncertain{monômes} \ill{}
\[
\xi_i = x_1^{(p_1)}\cdots x_r^{(p_r)}
\]
\[
u_\alpha\Bigl(\bigl(x_1^{(p_1)}\cdots x_r^{(p_r)}\bigr)^{(q)}\Bigr) \overset{?}{=}
u_\alpha\bigl(x_1^{(p_1)}\cdots x_r^{(p_r)}\bigr)^{(q)}
\]
où
\[
\bigl(x_1^{(p_1)}\cdots x_r^{(p_r)}\bigr)^{(q)} =
\bigl(x_1^{(p_1)}\cdots x_{r-1}^{(p_{r-1})}\bigr)^{q}\bigl(x_r^{(p_r)}\bigr)^{q}
\]
\[
u_\alpha\Bigl(\bigl(x_p^{(p)}\bigr)^{(q)}\Bigr) = u\bigl(\alpha_{pq}x^{(pq)}\bigr)
\]
\struck{$u(\ldots)=$}
\[
u\bigl((x^{(p)})^{(q)}\bigr) = u(x^{(p)})^{(q)}, \qquad F(\gamma^p_\alpha)(x) = x^{(p)}
\]

$A \xrightarrow{u} B$, $u$ \uncertain{compatible} avec \ill{}~? \quad $\check J \to \check K$, \quad $u(\ldots)$
$(x_i)_{i\in I}$ \uncertain{engendrant} par les $x_i$ comme \uncertain{idéal} à p.d., i.e.\ comme \uncertain{module} \uncertain{engendré} par les
\[
x_i^{(p)} \qquad (p\ge1)
\]
Si $u(x_i^{(p)}) = u(x_i)^{(p)}$ ok.
\[
u\Bigl(\bigl(\textstyle\sum\lambda_ix_i^{(p_i)}\bigr)^{(q)}\Bigr) \overset{?}{=}
\Bigl(u\bigl(\textstyle\sum\lambda_ix_i^{(p_i)}\bigr)\Bigr)^{(q)}
\qquad (\xi_i = \lambda_ix_i^{(p_i)})
\]
\[
u\Bigl(\textstyle\sum\xi_1^{(q_1)}\cdots\xi_r^{(q_r)}\Bigr) =
\textstyle\sum u(\xi_1)^{(q_1)}\cdots u(\xi_r)^{(q_r)},
\]
il suffit $u(\xi_i^{(q_i)}) = u(\xi_i)^{(q_i)}$.

OPS.\ $x = \lambda x_i^{(p)}$,
\[
u\bigl(x^{(q)}\bigr) = u(x)^{(q)}, \qquad
u\bigl(\lambda^q(x_i^{(p)})^{(q)}\bigr) = u(\lambda)^q\bigl(u(x_i)^{(p)}\bigr)^{(q)}
\]

\page{333}
\uncertain{Plus généralement}, si $A \supset J$ est un idéal, \uncertain{revient} à \uncertain{se donner} \ill{} $F(J) \to F(A)$ \uncertain{injectif} dans $F(J)$, à p.d.\ dans $F(A)$~?
$F(J) \xrightarrow{\deg p} F(J)$ oui, en \uncertain{transformant} $J \xrightarrow{\gamma^p} J$.
\struck{\ill{}} \uncertain{Il faut} \uncertain{donner la} \uncertain{valeur}
\struck{$(x+y)$}
\[
x,y\in F(J) \Longrightarrow (x+y)^{(p)} = \textstyle\sum x^{(i)}y^{(j)}\;? \quad (\deg p)
\]
$x,y\in$ \; $u$
\[
\mathcal{P} \xrightarrow{\;F\;} \mathcal{P}'
\]
\uncertain{commutant} aux produits \uncertain{tensoriels}, \uncertain{commutant} aux produits et aux \uncertain{degrés partiels}, \qquad bidegrés 1,1.

\struck{Lemme} \add{il} : $A$, $J$, \quad $J \to A$, \quad $A\times A \to A$, \quad $A\otimes J \to J$ ;

$A$ \uncertain{$k$-algèbre assoc.\ commutative} avec 1,

$J$ \uncertain{module assoc.\ comm.} sans 1,

$A$ opérant sur $J$, et $(\lambda x)y = x(\lambda y) = \lambda xy$,

$J \xrightarrow{\;i\;} A$ hom.\ de $A$-\uncertain{modules} ($i(xy) = i(x)i(y)$, $i(\lambda x) = \lambda i(x)$),

$k \to \mathrm{Hom}(\mathcal{O},A)$ hom.\ d'anneaux.
\[
\gamma^p : J \to J \quad (p\ge1) \;\; \text{de degré } p, \qquad \gamma^0 : J \to A
\]
\[
\begin{aligned}
\gamma^0 &= 1 \\
\gamma^1 &= \mathrm{id}_J \\
\gamma^p(x+y) &= \textstyle\sum \gamma^ix\,\gamma^jy \\
\gamma^p(\lambda x) &= \lambda^p\gamma^p(x) \\
\gamma^p(x)\gamma^q(x) &= c_{pq}\gamma^{p+q}(x) \\
\gamma^p(\gamma^q(x)) &= \frac{(pq)!}{(p!)^q\,q!}\,\gamma^{pq}(x)
\end{aligned}
\]
\note{après $\gamma^p : J\to J$, un mot barré illisible.}
\uncertain{Transformation naturelle}
\[
\Gamma^*F(X) \longrightarrow F\,\Gamma^*(X)
\]

\page{334}
\[
\Gamma^pF(X) \xrightarrow{\;u^X_p\;} F(\Gamma^p(X))
\]
\[
x\in F(X) \longmapsto x^{(p)} = u^X_p(\gamma^p(x))
\]
\[
u^X_* : \Gamma^*F(X) \longrightarrow F(\Gamma^*(X))
\]
hom.\ d'\uncertain{algèbres}
\note{sous $\Gamma^*(X)$ : « \uncertain{algèbre} \uncertain{car} $F$ compatible à $\otimes$ \uncertain{donc} \uncertain{transforme} \uncertain{alg.\ en alg.} ».}
\[
\begin{array}{c}
F(P)\otimes F(Q) \\
\downarrow \\
F(P\otimes Q) \to F(R)
\end{array}
\]
\uncertain{calcul} \struck{\ill{}} \add{\uncertain{implique}} \uncertain{pour}
\[
(x+y)^{(p)} = \sum_{i+j=p} x^{(i)}y^{(j)}
\]
\[
u_*\bigl(\gamma_p(x+y)\bigr) = \sum_i u_*\bigl(\gamma_i(x)\gamma_j(y)\bigr)
= \sum u(\gamma_i(x))\,u(\gamma_j(y)) = \sum x^{(i)}y^{(j)}
\]
\[
(\lambda x)^{(p)} = \lambda^px^{(p)} :
\]
\[
u^X_p\gamma^p(\lambda x) = u^X_p\bigl(\lambda^p\gamma^p(x)\bigr) = \lambda^pu^X_p(\gamma^p(x)) = \lambda^px^{(p)}
\]
\[
x^{(p)}x^{(q)} = c_{pq}\,x^{(p+q)}\;? :
\]
\[
u_*\gamma^p(x)\,u_*\gamma^q(x) = u_*\bigl(\gamma^p(x)\gamma^q(x)\bigr) = c_{pq}\,u\gamma^{p+q}(x) = c_{pq}\,x^{(p+q)}
\]

À l'encre bleue :
\[
\mathbb{1} \longrightarrow F(\mathbb{1}), \qquad
\Gamma^*F(X) \longrightarrow F\,\Gamma^*(X), \qquad
F(X) \xrightarrow{\;\mathrm{id}\;} F(X)
\]
\uncertain{multiplicatif}, \ill{} \uncertain{foncteur} \ill{} p.d.~?

$F\,\Gamma^*(X)$ \uncertain{est-il} \uncertain{à} p.d.~?

\begin{tikzcd}[column sep=small, row sep=small, nodes={font=\scriptsize}]
\Gamma^pF(X) \arrow[r, "u^{F(X)}_p"] \arrow[d, "\gamma^q_{\Gamma^pF(X)}"'] & F\Gamma^p(X) \arrow[d, "F(\gamma^q_{\Gamma^p(X)})"] \\
\Gamma^q\Gamma^pF(X) \arrow[d] & F(\Gamma^q\Gamma^p(X)) \\
\Gamma^{qp}F(X) & {}
\end{tikzcd}

\section{Gribouillis}
\note{le mot « Gribouillis », souligné, est seul sur la p.~335, de sa main ; il sert de titre aux feuillets suivants.}

\page{336}
\[
\begin{array}{ccccc}
 & & \Lambda^n\Psi_* & \longrightarrow & \Lambda^{n+1}\Phi_* = C^n_* \\
 & & \uparrow & & \uparrow \\
\Lambda^{n-1}\Psi_* & \xrightarrow{\;T_\wedge\;} & \Phi_*\otimes\Lambda^{n-1}\Psi_* & \longrightarrow & \Lambda^n\Phi_* = C^{n-1}_* \\
\uparrow & & \uparrow & & \uparrow \\
\Phi_*\otimes\Lambda^{n-2}\Psi_* & \longrightarrow & \Gamma^2\Phi_*\otimes\Lambda^{n-2}\Psi_* & \longrightarrow & \Lambda^{n-1}\Phi_* = C^{n-2}_*
\end{array}
\]
\note{la flèche $T_\wedge$ porte en dessous « \uncertain{mono} » ; sous la deuxième ligne on lit $(\Lambda^n\Phi_*)^{\binom n2}$, sous la troisième une expression $(\Lambda^n\Phi_*+\Phi\Lambda^{\ldots}\Phi_*)^{\ldots}$ dont les exposants sont illisibles.}
\[
(N+1)\binom{N}{n-1} = \frac{(N+1)N\cdots(N-n+2)}{(n-1)!} = n\binom{N+1}{n}
\]
\[
\frac{(N+1)(N+2)}{2}\binom{N}{n-2} = \frac{(N+2)(N+1)N\cdots(N-n+3)}{n!}\,\frac{n(n-1)}{2}
= \binom{N+2}{n}\frac{n(n-1)}{2}
\]
\note{au-dessus du numérateur : « $n$ facteurs ».}
\[
\mathrm{rg}\bigl(\Gamma^p\Phi_{(N)}\otimes\Lambda^q\Psi_{(N)}\bigr)
= \frac{(N+1)(N+2)\cdots(N+p)}{p!}\cdot\frac{N(N-1)\cdots(N-q+1)}{q!}
\]
\[
= \frac{(N+p)\cdots(N-q+1)}{(p+q)!}\,\frac{(p+q)!}{p!\,q!} = c_{pq}\binom{N+p}{p+q}
\]
\[
(\tau P)(N) = P(N+1)
\]
\[
\begin{aligned}
\tau\varphi_r &= \varphi_r + \varphi_{r-1} \\
\tau^2\varphi_r &= \tau\varphi_r + \tau\varphi_{r-1} = \varphi_r + \varphi_{r-1} + \varphi_{r-1} + \varphi_{r-2}
= \varphi_r + 2\varphi_{r-1} + \varphi_{r-2} \\
\tau^3\varphi_r &= \varphi_r + \varphi_{r-1} + 2\varphi_{r-1} + 2\varphi_{r-2} + \varphi_{r-2} + \varphi_{r-3} \\
&= \varphi_r + 3\varphi_{r-1} + 3\varphi_{r-2} + \varphi_{r-3} \\
\tau^p\varphi_r &= \varphi_r + p\,\varphi_{r-1} + \binom p2\varphi_{r-2} + \cdots + \varphi_{r-p} \\
\tau^{p-1}\varphi_{p+q} &= \varphi_{p+q} + \binom{p-1}{1}\varphi_{p+q-1} + \cdots + \varphi_{q+1}
\end{aligned}
\]
\[
c_{p,q}\Bigl(\Lambda^{p+q}\Phi_* + (p-1)\Lambda^{p+q-1}\Phi_* + \binom{p-1}{2}\Lambda^{p+q-2}\Phi_* + \cdots + \Lambda^{q+1}\Phi_*\Bigr)
\]
À droite :
\uncertain{Mais}
\[
\varphi_{p+q}(N+p) = \varphi_{p+q}(N+p-1) + \varphi_{p+q-1}(N+p-1)
\]
\[
= \varphi_{p+q}(N+p-1) + 2\varphi\ldots
\]
\[
= c_{pq}\,\tau^{p-1}\varphi_{p+q}(N+1)
\]
\note{la deuxième ligne est inachevée ; quelques essais barrés à côté ne sont pas lisibles.}

\page{337}
\note{feuillet écrit en travers, à l'encre bleue ; formules disposées autour de deux hexagones de flèches. On les transcrit dans l'ordre de lecture.}

$\mathrm{Hom}^!(\Delta^m,\Delta^n) =$ \uncertain{ens.} des applications croissantes $\Delta^m \to \Delta^n$ \struck{$[0,n]$} dont l'image \uncertain{ne soit pas contenue} \ill{} $\Delta^n$ (\ill{}) \uncertain{ou} $[0,n-1]$.

\drawing{hexagone de flèches reliant $\mathbb{C}^*_*$, $\mathbb{C}^\bullet_*$, $\mathbb{C}^*_\bullet$ et $\mathbb{C}^\bullet_\bullet$}
\[
\begin{cases}
\mathbb{C}^*_{[n]} = \mathbb{C}^*(\Delta^n_*) = k_*(\Delta^n_*) = k^{\Delta^n_*} \\
\mathbb{C}^{[m]}_* = p_{m*}(k) = \bigl(n\mapsto k^{\mathrm{Hom}_\Delta(\Delta^m,\Delta^n)}\bigr)
\end{cases}
\]
\[
\mathbb{C}^\bullet_* = \Lambda^{\bullet+1}\Phi_* \quad
\begin{cases}
\mathbb{C}^\bullet_{[n]} = \mathbb{C}^\bullet(\Delta^n) = \Lambda^{\bullet+1}k^{\Delta^n} \\
\mathbb{C}^m_* = \Lambda^{m+1}\Phi_*
\end{cases}
\]
\[
\mathbb{C}^{(m)}_{[n]} = k^{\mathrm{Hom}_\Delta(\Delta^m,\Delta^n)}, \qquad
\mathbb{C}^{[m]}_n = k^{\mathrm{Hom}^!(\Delta^m,\Delta^n)}
\]
\[
\mathbb{C}^m_{[n]} = \Lambda^{m+1}k^{\Delta^n} = k^{\mathrm{Inj}_\Delta(\Delta^m,\Delta^n)}
\quad (\text{nul si } m>n,\; = k \text{ si } m=n)
\]
\note{des flèches entre ces objets portent « isomorphisme si $m=0$ », « fact.\ dir.\ can. », « nul si $n>m+1$ ; $= k$ si $n=m+1$ », « iso si $n\le m+1$ », « iso si $m\ge n$ » (ce dernier barré).}
\[
\mathbb{C}^m_n = \begin{cases} 0 & \text{si } n\neq m, m+1 \\ k & \text{si } n = m \text{ ou } n = m+1 \end{cases}
\]
alors $\mathbb{C}^m_{[m]} = \Lambda^{m+1}k^{m+1} \simeq k$, \; $\mathbb{C}^m_{[m+1]} = \Lambda^{m+1}k^{m+2} \simeq k^{m+2}$ et \ill{} l'injection \uncertain{diagonale} \ill{}.
\[
\mathbb{C}^m_\bullet \simeq k \text{ en degrés } m, m+1 \quad (\text{dans } (K_\bullet) ;\; \text{complexe de cochaînes } C^\bullet_\bullet)
\]
\[
\mathbb{C}^\bullet_m \simeq \begin{cases} k \text{ en degrés } m, m-1 & (m\ge1) \\ k \text{ en degré } 0 & \text{si } m = 0 \end{cases}
\]
\[
\mathrm{Hom}_*(X_*,\mathbb{C}^{[i]}_*) \simeq \mathbb{C}^{[i]}(X_*) \simeq k^{X_{[i]}} = \mathrm{Hom}_{\mathrm{ens}}(X_{[i]},k)
\]
\[
\mathrm{Hom}_*(X_*,\mathbb{C}^i_*) \simeq \mathrm{Hom}(X_*,\Lambda^{i+1}\Phi_*) = \mathbb{C}^i(X_*)
\]
\uncertain{sous-module de} $\mathbb{C}^{[i]}$
\[
\mathbb{C}_*(X_*) \overset{\mathrm{def}}{=} k(X_*) = k^{(X_*)} = \bigl(n\mapsto k(X_{[n]})\bigr) = k^{(X_{[n]})}
\]
\[
\mathbb{C}^*(X_*) \simeq \mathbb{C}_*(X_*)^\vee = \bigl(n\mapsto (k^{(X_{[n]})})^\vee = k^{X_{[n]}}\bigr)
\]
donc
\[
\mathbb{C}_\bullet(X_*) \subset \mathbb{C}_*(X_*),\;\ldots \qquad \mathbb{C}^\bullet(X_*) \simeq \mathbb{C}_\bullet(X_*)^\vee,
\qquad \mathbb{C}_\bullet(X_*) \simeq \mathrm{ND}(\mathbb{C}_*(X_*))
\]
\note{sous la dernière inclusion : « fact.\ dir.\ can. ».}

\drawing{second hexagone, dual du premier}
\[
\mathbb{C}_*{}^* = \mathbb{C}^*_*{}^\vee, \qquad
\mathbb{C}_\bullet{}^* = \mathbb{C}^\bullet_*{}^\vee = (\Lambda^{\bullet+1}\Phi_*)^\vee, \qquad
\mathbb{C}_*{}^\bullet = \mathbb{C}^*_\bullet{}^\vee, \qquad
\mathbb{C}_\bullet{}^\bullet = \mathbb{C}^\bullet_\bullet{}^\vee
\]
($\mathbb{C}_m^\bullet = k$ en deg.\ $m, m+1$) (dans $(K_\bullet)$ ; $\mathbb{C}_\bullet$ complexe de \uncertain{chaînes} \ill{})
\[
\mathbb{C}_\bullet{}^n = \begin{cases} k \text{ en degrés } n, n-1 & (n\ge1) \\ k \text{ en degré } 0 & (n=0) \end{cases}
\]

\page{338}
À tout $X_* \in \mathrm{Ob}\,\hat\Delta = \mathrm{Ob}\,\mathrm{Ss}$, on associe le foncteur
\[
F_{X_*} : M_* \longmapsto \mathrm{Hom}_{\hat\Delta}(X_*,M_*) \simeq \mathrm{Hom}_{\hat\Delta_{k\text{-}\mathrm{lin}}}(k[X_*],M_*) \in \mathrm{Ab}_k
\]
\[
k[X_*] = \lambda_* = C_*(X_*), \qquad n \longmapsto C_{[n]}(X_*) = k[X_{[n]}] = k^{(X_n)}
\]
\[
\Bigl(F_{X_*} \in \underline{\mathrm{Hom}}_{k\text{-}\mathrm{lin}}(\mathrm{Ab}_*,\mathrm{Ab}_k)\Bigr) \simeq G_{(\lambda_*)}(M_*), \qquad G_{\lambda_\bullet}(M_\bullet)
\]
La donnée d'un tel foncteur $F$ \uncertain{$k$-lin.\ exact à gauche}, \ill{} \uncertain{on a}
\[
F(\mathbb{C}^\bullet_*) \in \mathrm{Ab}^\bullet_k, \qquad
F(\mathbb{C}^\bullet_*) = \bigl(i\mapsto F(\mathbb{C}^i_*)\bigr) \in \mathrm{Ab}_k
\]
On a
\[
F_{X_*}(\mathbb{C}^i_*) \simeq \mathbb{C}^i(X_*) = k^{X_{[i]}} \;\simeq\; \mathbb{C}_i(X_*)^\vee
\]
(cochaînes \uncertain{non dég.})
\marginal{entouré : plus gén., $G_{\lambda_*}(\mathbb{C}^i_*) \simeq \lambda_i^\vee$, \; $= G_{\lambda_\bullet}(\mathbb{C}^\bullet)$}
\[
F_{X_*}(\mathbb{C}^\bullet_*) \simeq \mathbb{C}^\bullet(X_*)
\]
$=$ \uncertain{complexe des cochaînes non dégénérées} sur $X_*$
Ex.\ $X_* = \Delta^n_*$, \; $F_{\Delta^n_*}(M_*) \simeq M_{[n]}$,
\[
\mathbb{C}^\bullet(\Delta^n) \simeq (\mathbb{C}^i_{[n]})_i = \Lambda^{*+1}k^{\Delta^n} = \bigl(m\mapsto k^{\mathrm{Hom}^!(\Delta^m,\Delta^n)}\bigr)
\]
On s'intéressera, pour $X_*$ fixe, à
\[
\mathrm{Hom}(F_{X_*},F_{\Delta^n_*}) \simeq \mathrm{Hom}_\bullet\bigl(\mathbb{C}^\bullet(X_*),\mathbb{C}^\bullet(\Delta^n)\bigr),
\qquad \lambda^\vee_\bullet, \quad \Lambda^{*+1}k^{\Delta^n} = \mathbb{C}^\bullet_{[n]}
\]
\[
\simeq \mathrm{Hom}(\mathbb{C}^\bullet_{[n]},\lambda^{\vee\vee}_\bullet)
\;\overset{!}{\simeq}\; (\lambda^{\vee\vee}_\bullet)_{[n]} = \lambda_{[n]}^{\vee\vee} = k[X_{[n]}]^{\vee\vee}
\]
\struck{$\mathrm{Hom}(G_{\lambda_*},G_{\mathbb{C}_*(\Delta^n)})$} \quad $u \mapsto F(u)$, \quad $\lambda_* = \mathbb{C}_*[X]$
\[
k[X_n] \longleftarrow \mathrm{Hom}(\Delta^n,X_*) \simeq X_n
\]
\struck{$K = \mathrm{Ab}_{k*}$}
\uncertain{Avons} \ill{} \ill{} \uncertain{diagramme commutatif}

\begin{tikzcd}[column sep=small, row sep=small, nodes={font=\scriptsize}]
X_{[n]} \simeq \mathrm{Hom}(\Delta^n_*,X_*) \arrow[d] \arrow[drr] & {} & {} \\
k[X_{[n]}] \arrow[d, "\text{iso si } X_n \text{ fini}"'] & {} & {} \\
k[X_{[n]}]^{\vee\vee} & {} & \mathrm{Hom}(F_{X_*},F_{\Delta^n_*}) \arrow[ll, "\sim"']
\end{tikzcd}

\[
\mathrm{Hom}(F_{X_*},F_{\Delta^n_*}) \simeq \mathrm{Hom}_\bullet\bigl(\mathbb{C}^\bullet(X_*),\mathbb{C}^\bullet(\Delta^n_*)\bigr)
\]
(\struck{$\mathrm{Hom}(\ldots)$})
\[
\Bigl(\mathrm{Hom}_\bullet\bigl(\Lambda^{*+1}k^{\Delta^n},\mathbb{C}^\bullet(X_*)^\vee\bigr)
\simeq \mathrm{Hom}\bigl(\mathbb{C}^\bullet(X_*),\Lambda^{*+1}k^{\Delta^n}\bigr)\Bigr)
\]
Faisons en particulier $X_* = \Delta^m_*$, on trouve donc
\[
\begin{array}{ccc}
\Delta^m_{[n]} = \mathrm{Hom}(\Delta^n,\Delta^m) & \simeq & \mathrm{Hom}(\Delta^n_*,\Delta^m_*) \\
\downarrow & & \downarrow \\
k(\Delta^m_{[n]}) & \xrightarrow{\;\sim\;} & \mathrm{Hom}(F_{\Delta^m_*},F_{\Delta^n_*}) \simeq
\end{array}
\]
\note{la formule s'interrompt au bas de la page ; elle reprend p.~339.}

\page{339}
\[
k^{\mathrm{Hom}(\Delta^m,\Delta^n)} \longleftarrow k^{\mathrm{Hom}(\Delta^m,\Delta^{n-1})}
\]
\[
C^\bullet_{[n]} = \Lambda^{\bullet+1}\Phi_{[n]} = \Lambda^{\bullet+1}k^{\Delta^n} \longleftrightarrow \lambda^\bullet
\]
\[
C_\bullet^{[n]} = \Lambda^{\bullet+1}\check\Phi_{[n]} = \Lambda^{\bullet+1}k^{\Delta^n} \longrightarrow \lambda_\bullet
\]
\[
C_*^{[n]} \longrightarrow \lambda_*, \qquad \mathbb{C}^*_{[n]} = \mathbb{C}^*(\Delta^n_*)
\]
\[
C_*^{[n]} = \mathbb{C}^*(\Delta^n_*)^\vee = \mathbb{C}_*(\Delta^n_*) = k[\Delta^n_*], \qquad
\mathbb{C}^\bullet_{[n]} \to K^\bullet, \quad \mathbb{C}^*_{[n]} \to K^*
\]
\[
\mathrm{Hom}_k(k[\Delta^n_*],\lambda_*) = \mathrm{Hom}(\Delta^n_*,\lambda_*) = \lambda_{[n]}
\]
\uncertain{prenons} $\lambda^\bullet = \mathbb{C}^\bullet(\Delta^m_*) \simeq \Lambda^{1+\bullet}k^{\Delta^m}$.
\[
\mathrm{Hom}^\bullet\bigl(\Lambda^{1+\bullet}k^{\Delta^m},\Lambda^{1+\bullet}k^{\Delta^n}\bigr) = \mathrm{Hom}^\bullet(C^\bullet_{[m]},\lambda^\bullet)
\]
\note{sous les deux arguments : « \uncertain{avec produit ext.} ».}
\[
\mathrm{Hom}_\bullet\bigl(\Lambda^{1+\bullet}k^{\Delta^m},\Lambda^{1+\bullet}k^{\Delta^n}\bigr) \simeq \mathrm{Hom}_\bullet(C_\bullet^{[m]},\lambda_\bullet)
\]
\note{sous les deux arguments : « \uncertain{avec prod.\ int.} ».}
\[
\simeq \quad \lambda_{[n]} = \mathbb{C}_{[n]}(\Delta^m_*) = \mathbb{C}^{[n]}(\Delta^m_*)^\vee
\]
\struck{$= \bigl(k^{\mathrm{Hom}(\Delta_n,\Delta_m)}\bigr)^\vee$}
\[
\simeq k^{(\mathrm{Hom}(\Delta_n,\Delta_m))}
\]
\[
k^{\mathrm{Hom}(\Delta_n,\Delta_m)} \longrightarrow \mathrm{Hom}_k\bigl((k^{\Delta_m},T_m),(k^{\Delta_n},T_n)\bigr)
\]
\uncertain{pas bijectif}~!
\[
\searrow\;?\quad \mathrm{Hom}\bigl(\Lambda^{\bullet\bullet}k^{\Delta_m},\Lambda^{\bullet\bullet}k^{\Delta_n}\bigr)
\]
\note{une longue flèche relie l'expression $\simeq k^{(\ldots)}$ au $\mathrm{Hom}_\bullet$ avec produit intérieur.}

\page{340}
\[
L_\bullet \longmapsto \mathrm{Hom}(\lambda_\bullet,L_\bullet)
\]
\[
\mathrm{Hom}(\lambda_\bullet,\mathbb{C}^i_\bullet) \simeq \lambda_i^\vee, \qquad
\mathrm{Hom}(\lambda_\bullet,\mathbb{C}^\bullet_\bullet) \simeq \lambda^\vee_\bullet
\]
\[
\mathrm{Hom}\bigl(\lambda^\vee_\bullet,\Lambda^{*+1}k^{\Delta_n}\bigr) \simeq \mathrm{Hom}_\bullet\bigl(\mathbb{C}^{\bullet\vee}_{[n]},\lambda^{\vee\vee}_\bullet\bigr),
\qquad \Lambda^{*+1}k^{\Delta_n} = \mathbb{C}^\bullet_n, \quad \mathbb{C}^{\bullet\vee}_{[n]} = C_\bullet^{[n]}, \quad \lambda^{\vee\vee}_\bullet = \lambda_\bullet
\]
\struck{\ill{}} \add{$k$-module} \uncertain{à simplicial} \uncertain{pour une variable}

\[
\begin{array}{ccc}
C_i{}^n = \Lambda^{n+1}k^n & & \\
\downarrow{\scriptstyle\partial^*_*} & & \\
C_i{}^n & \longrightarrow & \tilde\lambda_i \\
\| & & \downarrow{\scriptstyle d_i} \\
\Lambda^{i+1}k^{n+1} & & \\
\downarrow & & \\
C_{i-1}{}^n & \longrightarrow & \tilde\lambda_{i-1}
\end{array}
\]
\note{la seconde flèche verticale de gauche est marquée « produit intérieur par $T_n$ » ; entre les colonnes : « \uncertain{comm.} ».}

\struck{$\mathrm{Hom}(\lambda_\bullet,\mathbb{C}^i_n)$} \quad \struck{$\mathrm{Hom}_\bullet(\lambda_\bullet,\mathbb{C}^i_\bullet)$ $i$ variable} \quad \struck{$i \mapsto \mathbb{C}^i_\bullet$ \ill{} \uncertain{cosimplicial}}

\[
K_{n-1} \underset{k^*}{\overset{\partial^*}{\rightleftarrows}} K_n \;\longmapsto\;
\bigl(n\mapsto \mathrm{Hom}_\bullet(\mathbb{C}_\bullet^{[n]},\tilde\lambda_\bullet)\bigr) = \mathrm{Hom}_\bullet(C_\bullet{}^*,\tilde\lambda_\bullet),
\qquad
K^{n-1} \underset{k_*}{\overset{\partial_*}{\rightleftarrows}} K^n
\]
\struck{$\mathrm{Hom}(\lambda_\bullet,\ldots)$}

$n \mapsto C_\bullet^{[n]}$ est un objet cosimplicial de $\mathrm{Ab}_{k\bullet}$, qui \uncertain{transformé} par le foncteur \uncertain{contravariant} $\mathrm{Hom}_\bullet(-,\tilde\lambda_\bullet)$ \uncertain{à valeurs dans} $\mathrm{Ab}_k^\bullet$, \uncertain{donne un} \uncertain{objet simplicial} de $\mathrm{Ab}^\bullet_k$ \uncertain{(ou complexe de} $\mathrm{Ab}_k$) \uncertain{dont} \ill{} \uncertain{le complexe associé est} AD, \ill{} \uncertain{on peut revenir} \ill{} \uncertain{que} \ill{} \uncertain{formule le} AD \uncertain{en partant} de $\mathbb{C}_\bullet{}^*$ i.e.\ $C_\bullet{}^\bullet$.

\[
(\mathbb{C}^i{}_*)_i \simeq C^\bullet_*, \qquad C^i_{[n]} = \Lambda^{i+1}k^{n+1}
\]
\note{à côté : « \uncertain{produit ext.\ par les vecteurs de base} de $k^{n+1}$ » et, au-dessus, « \uncertain{pour degré} $i$ ».}
\[
(\mathbb{C}_i{}^*)_i = \mathbb{C}_\bullet{}^* = (\mathbb{C}^\bullet_*)^\vee, \qquad
\mathbb{C}_i{}^{[n]} \simeq C^i_{[n]}{}^\vee \simeq \Lambda^{i+1}k^{n+1}
\]
(\uncertain{produit int.\ par les} \ill{} \uncertain{pour} \ill{})
\[
(\mathbb{C}^i_\bullet)_i = C^\bullet_\bullet, \qquad
C^i_j = \begin{cases} 0 & \text{si } j\neq i, i+1 \\ k & \text{si } j = i, i+1 \end{cases}
\]
\[
(\mathbb{C}_i{}^\bullet)_i = C_\bullet{}^\bullet = C^\bullet_\bullet{}^\vee, \qquad
C_i{}^j = C^i_j{}^\vee \simeq \begin{cases} 0 & \text{si } j\neq i, i+1 \\ k & \text{si } j\neq i, i+1 \end{cases}
\]
\note{sic : la seconde condition répète « $j\neq i,i+1$ » pour « $j = i,i+1$ ». Rien n'indique que l'argument se poursuive au-delà du lot.}

\end{document}
